% ============================================================
%  Physical-Contact (Touch) Detection — Methods Explored & Results
%  Smart Thermal System for Patient Safety Monitoring
%  Waveshare 26984 (80x62) Thermal Sensor Dataset
%  Afeka College of Engineering, Tel-Aviv
% ============================================================
\documentclass[11pt, a4paper]{article}

\usepackage[a4paper, margin=2.4cm]{geometry}
\usepackage[T1]{fontenc}
\usepackage{lmodern}
\usepackage{booktabs}
\usepackage{array}
\usepackage{multirow}
\usepackage{amsmath}
\usepackage[table]{xcolor}
\usepackage{hyperref}
\usepackage{parskip}
\usepackage{enumitem}
\usepackage{float}

\providecommand{\arraybackslash}{\let\\\tabularnewline}
\makeatletter
\providecommand\do@row@strut{}
\let\do@row@strut\relax
\makeatother

\definecolor{headerblue}{RGB}{31,73,125}
\definecolor{rowhi}{RGB}{223,240,216}
\hypersetup{colorlinks=true, linkcolor=headerblue, urlcolor=headerblue}

\title{%
    {\large Afeka College of Engineering, Tel-Aviv}\\[0.3em]
    {\large Smart Thermal System for Patient Safety Monitoring}\\[0.7em]
    {\LARGE\bfseries Physical-Contact (Touch) Detection}\\[0.2em]
    {\large Methods Explored and Results}\\[0.3em]
    {\normalsize Waveshare 26984 \quad 80$\times$62\,px \quad 3 synchronised views}
}
\author{Guy Chen \and Yaniv Blau \and Roy Lieberman\\[0.2em]
    \small Supervisor: Or Zilberberg \qquad Advisor: Dr.\@ Oshrit Hoffer}
\date{June 2026}

\begin{document}
\maketitle

% ------------------------------------------------------------
\section*{1.\quad Objective and setup}
The touch module must flag \emph{physical contact between patients} from three
synchronised thermal cameras (80$\times$62\,px, 8\,Hz), with no optical imagery.
Contact is annotated per frame in \texttt{contact\_labels.csv}:
\textbf{188 positive frames across only four scenes} (\texttt{2ppl\_fight},
\texttt{2ppl\_hug}, \texttt{3pp\_surprise}, \texttt{3pplhedroncolider}); all
other scenes are contact-negative, including hard negatives where people are
close but not touching (\texttt{2ppl\_walk}, \texttt{3ppl\_dance}). All methods
are scored on a common per-session timeline split (test = 695 frames, 41
positive). The primary metric is \textbf{F1}; recall is the safety priority,
false-alarm rate guards against alert fatigue.

% ------------------------------------------------------------
\section*{2.\quad Approaches tried and results}

We evaluated three detector families from the design (multi-view geometric,
spatiotemporal graph network, volumetric appearance network), then developed a
homography-free pipeline that progressively outperformed them. Table~1 gives the
headline F1 on the common test split.

\begin{table}[H]
\centering
\caption{Touch detectors on the common test split (single split; see §4 for the
cross-validated picture).}
\renewcommand{\arraystretch}{1.25}
\begin{tabular}{l l c}
\toprule
\textbf{Family / method} & \textbf{Idea} & \textbf{F1} \\
\midrule
Geometric multi-view      & boxes $\to$ homography $\to$ floor distance & 26.8\% \\
MV-STGCN                  & graph net on projected actor tracks         & 25.2\% \\
Thermo-X3D                & 3-D CNN on raw thermal volume (no boxes)    & 41.7\% \\
\midrule
\multicolumn{3}{l}{\textit{Homography-free image-plane pipeline (developed here):}}\\
\;Box-proximity voting    & per-camera box gap + 3-camera vote          & 28.6\% \\
\;+ merge-aware / routing  & merged-blob counts; people-count routing    & 35.5\% \\
\;+ blob-merge            & contact = two people in one warm blob       & 44.6\% \\
\;+ temporal smoothing    & morphological open/close over time          & 53.8\% \\
\;+ two-body merge        & require exactly-2-body merge in crowds       & 55.0\% \\
\;\textbf{+ raw detector (final)} & \textbf{detector on raw, blobs on residual} & \textbf{60.8\%} \\
\bottomrule
\end{tabular}
\end{table}

\noindent The final detector (``config-D'') chains: MobileNet-SSD person
detection $\to$ warm-blob segmentation on the background-subtracted residual
$\to$ flag a frame when exactly two people occupy one connected warm component
$\to$ temporal smoothing. It uses \textbf{no homography} and needs no contact
training labels (thresholds only). It surpasses the deep Thermo-X3D baseline.

% ------------------------------------------------------------
\section*{3.\quad Key findings}
\begin{itemize}[leftmargin=1.4em, itemsep=2pt]
\item \textbf{Homography is the wrong lever here.} With no surveyed floor
      points, the homography was self-calibrated from a walking subject; it was
      accurate enough to localise but too noisy to fuse cleanly, so the
      geometric and graph methods underperformed (F1 $\approx$25--27\%).
\item \textbf{The thermal ``merge'' signal is decisive.} Two people in contact
      fuse into one connected warm region; testing for a \emph{two-body} merged
      blob (rejecting three-body clusters) was the single biggest gain and is
      simple and explainable.
\item \textbf{Preprocessing should differ by task.} Background subtraction
      (Tateno) is \emph{essential} for blob segmentation but \emph{hurts} the
      person detector --- training/running the detector on \emph{raw} frames
      while segmenting blobs on the \emph{residual} added $\approx$6 F1 points.
\item \textbf{Temporal smoothing helps} (+9 points): contacts persist over time,
      so removing one-frame false alarms and filling one-frame dropouts is
      effective.
\item \textbf{Added complexity did not help.} A learned feature-combiner,
      motion/optical-flow features, and joint hyper-parameter tuning all matched
      or \emph{underperformed} the simple rule --- consistent symptoms of too
      few training examples.
\end{itemize}

% ------------------------------------------------------------
\section*{4.\quad Honest evaluation: the data limit}
A single test split overstates performance. Under 4-fold cross-validation
(every frame tested once, no detector leakage) the final detector scores
\textbf{F1 $\approx$44\% (fold mean 37\%, std $\pm$24)}. A bootstrap gives a
scene-level 95\% confidence interval of \textbf{[12\%, 66\%]} --- because
positives come from only four scenes, the estimate spans almost the whole range.
The 60.8\% above is a favourable split, not a reliable figure.

\paragraph{Can more data fix it?} We quantified the requirement and tested a
shortcut:
\begin{itemize}[leftmargin=1.4em, itemsep=2pt]
\item To measure F1 to $\pm$2.5\% (so 65\% is distinguishable from 60\%) needs
      on the order of \textbf{$10\times$ more positive frames and $3$--$4\times$
      more distinct contact scenes}.
\item \textbf{Re-using the older MLX90640 touch data} (real contact, but
      24$\times$32, upsampled) as pretraining \emph{reduced} F1 by 6 points ---
      the low resolution and sensor gap hurt more than the extra events helped.
      Synthetic compositing was judged unviable (multi-view consistency).
\end{itemize}

% ------------------------------------------------------------
\section*{5.\quad Conclusion and recommendation}
A practical, explainable, homography-free touch detector exists --- MobileNet-SSD
+ residual blob-merge + temporal smoothing (``config-D'') --- which on a matched
test split reaches 60.8\% F1 and beats the deep baseline. However, the
cross-validated estimate ($\approx$44\%, with a very wide confidence interval) shows
the system is \textbf{limited by data, not by algorithm}: 188 positive frames in
four scenes cannot reliably train richer models or measure performance.

\textbf{Recommendation:} adopt config-D as the working method and prioritise
\textbf{collecting more contact recordings across more distinct scenarios}
(target: several times the current positive count and contact scenes). This is
the single change expected to both raise and stabilise performance; further
algorithmic tuning on the current data is within measurement noise.

\end{document}
